\scanpage{57}
\mspage{II. 5}
% lines 1-4
alors, \deleted{\unsure{$k$}} $U$ étant normal, $\pi_1(U,\unsure{\xi|\zeta})\mathrel{\text{---}}\pi_1(X,\xi)$ est surjectif, donc il \deleted{est ramené au} suffit de prouver l'assertion pour $U,\xi$ et $A|U$. Cela nous permet de supposer $X$ affine, donc de t.f. sur $k$.

% lines 5-9
On sait alors que $X,A$ provient de $S_0=\operatorname{Spec}k_0$, $k_0$ un sous-corps de t.f. de $k$, $X_0$ schéma de t.f. sur $S_0$, $A_0$ schéma abélien sur $X_0$. Alors, \deleted{\unsure{$\xi$ définit un pt}} $X\mathrel{\text{---}}X_0$ étant fid. plat et $X$ normal, $X_0$ est normal.

\diagram{\begin{tikzcd}[ampersand replacement=\&]
A_0\arrow[r,dash]\arrow[d,dash] \& A\arrow[d,dash]\\
X_0\arrow[r,dash]\arrow[d,dash] \& X\arrow[d,dash]\\
S_0\arrow[r,dash] \& S
\end{tikzcd}}

% lines 10-13
De plus, $\xi$ définit un pt géométrique de $X_0$, donc on peut considérer \deleted{l'hom. de} la suite exacte du type (1.1.1)
\[
1\mathrel{\text{---}}\pi_1(\overline X_0,\xi)\mathrel{\text{---}}\pi_1(X_0,\xi)\mathrel{\text{---}}\pi_{\unsure{0|1}}(S_0,\xi)\longrightarrow\unsure{0}.
\]
% lines 14-17
\deleted{De plus, on \unsure{désigne le}} $\overline k_0$ est la clôture séparable de $k_0$ dans $k$, on a $\overline X_0\simeq X_0\otimes_{k}\overline k_0$, donc $\unsure{X_0|X}\simeq\overline X_0\otimes_{\overline k_0}k$. On a un hom \deleted{\unsure{de}}
\[
\pi_1(X,\xi)\longrightarrow\pi_1(\overline X_0,\xi)
\tag*{(1.5.1)}
\]
qui est en tous cas surjectif \deleted{\unsure{$X_0$ et}} (réf \quad).

% lines 19-25
On va considérer la représentation naturelle de $\pi_1(X,\xi)$ dans $T_\ell(A_{0\xi})$ \inserted{$=T_\ell(A_\xi)$}, et la représentation induite via $\pi_1(X,\xi)\longrightarrow\pi_1(\overline X_0,\xi)\hookrightarrow\pi_1(X,\xi)$ n'est autre que celle qu'il s'agit d'étudier. \deleted{\unsure{À}} Par la surjectivité de (1.5.1), dire qu'elle \deleted{\unsure{est que}} \inserted{\unsure{est} ess. résoluble (resp.} \inserted{ess\textsuperscript{t}} unipotente) signifie que celle induite sur $\pi_1(\overline X_0,\xi)$ l'est. \deleted{Donc} On conclut donc grâce à 1.3.
% archive sheet number: 56
